% ===================================================================
%  TABELL Nr 9 — Berget. — Badorten. — Skogen. — Jakten.
%  Traduction 2026 d'apres texto/io, controlee sur texto/fr et
%  texto/es. Consigne : voir texto/sv/10-tabell-01.tex.
% ===================================================================

%%K t09-apar-1 apar
\VUsaut{4.0mm}
\VUtitre{15.0pt}{60}{TABELL Nr 9}
\VUsaut{3.0mm}

%%K t09-tit-1 sub
\VUcentre{12.0pt}{0}{\textit{Första scenen.}}
\VUsaut{3.0mm}

%%K t09-tit-2 sub
\VUpk{11.4pt}{40}{Berget. --- Badorten.}
\VUsaut{3.0mm}

%%K t09-c1-01-1 p
1. --- Sedan åtta dagar är jag i Prats. Jag kan inte uttrycka hela den beundran jag kände när
jag stod framför Pyreneernas underbara \VUgras{bergskedja} \textsuperscript{(1)} vars
\VUgras{kam} \textsuperscript{(2)} tecknar sig mot den blå himlen som en jättelik festong.

%%K t09-c1-02-1 p
2. --- Jag föreställde mig inte, att Pyreneernas \VUgras{toppar} \textsuperscript{(3)} når en
så stor \VUgras{höjd}, och jag stod häpen inför de praktfulla \VUgras{glaciärerna}
\textsuperscript{(5)} och de snöiga
\VUgras{bergspetsarna} \textsuperscript{(4)} på vilka så förskräckliga
\VUgras{laviner} rullar.

%%K t09-tit-3 sub
\VUsaut{3.0mm}
\VUpk{10.2pt}{40}{Bergslandskap.}
\VUsaut{3.0mm}

%%K t09-c1-03-1 p
3. --- Redan dagen efter min ankomst gick jag till den gamla slocknade
\VUgras{vulkanen} \textsuperscript{(6)} som ligger vid Prats. På
\VUgras{kraterns} torra och kala kanter ser man ännu väl spåren av
\VUgras{lavans} framfart som flöt, för flera århundraden sedan, längs
\VUgras{bergssidan} \textsuperscript{(7)}. Det lyckades mig, på en av
\VUgras{sluttningarna}, att finna några stycken av den lavan.

%%K t09-c1-04-1 p
4. --- Prats ligger vid gränsen, och \VUgras{fästningen} \textsuperscript{(8)} som försvarar
\VUgras{bergspasset} \textsuperscript{(27)} som skiljer Frankrike från Spanien påminner helt
om de gamla slotten vid Rhens strand som från toppen av sin klippa tycks lura på ett
\VUgras{byte}.

%%K t09-c1-05-1 p
5. --- En enorm \VUgras{örn} \textsuperscript{(9)} som har sitt bo inte långt från
\VUgras{fortet} \textsuperscript{(8)} drar ofta till sig min uppmärksamhet. Den
\VUgras{rovfågeln}, vars \VUgras{näbb} är stark, bär ofta till sina ungar ett lamm eller en
killing som den håller med sina \VUgras{klor}. Så stor är dess
\VUgras{vingars} vidd att den länge kan sväva med sin tunga börda.

%%K t09-tit-4 sub
\VUsaut{3.0mm}
\VUpk{10.2pt}{40}{Bergsvandraren.}
\VUsaut{3.0mm}

%%K t09-c1-06-1 p
6. --- Där är den behagligaste promenaden utflykten till vattenfallet \textsuperscript{(10)}
vid « Cau ». \VUgras{Vattenfallet} störtar i virvlar ner och försvinner sedan som en vacker
\VUgras{bäck} i en \VUgras{granskog} \textsuperscript{(11)} väster om staden. Vi steg upp
genom den skogen ända till
\VUgras{det meteorologiska observatoriet} \textsuperscript{(12)} och från
\VUgras{utsiktsplatsens} topp såg vi hela traktens \VUgras{panorama}.
\VUgras{Landskapet} är underbart, och \VUgras{naturen} tycks ödsla på landet alla
de skatter den förfogar över.

%%K t09-c1-07-1 p
7. --- För att stiga ner använde vi \VUgras{bergbanan} \textsuperscript{(13)} och vi stannade
vid en liten \VUgras{station} bredvid \VUgras{ruinerna} \textsuperscript{(14)} av ett gammalt
\VUgras{feodalslott} som fordom beboddes av herrarna av Prats.

%%K t09-c1-08-1 p
8. --- Stackars slott! Vad tänker det väl när det nedanför ser den koketta
\VUgras{badorten} \textsuperscript{(15)} som nu är en \VUgras{modern kurort}?

%%K t09-c1-08-2 p
\VUgras{Klimatet} är så behagligt, \VUgras{mineralvattnet} så verksamt, att den
\VUgras{kur} som man använder i \VUgras{sanatoriet} \textsuperscript{(17)} är ett
sant nöje.

%%K t09-tit-5 sub
\VUsaut{3.0mm}
\VUpk{10.2pt}{40}{Behaglig kurstad.}
\VUsaut{3.0mm}

%%K t09-c1-09-1 p
9. --- För övrigt har man gjort allt för att göra vistelsen behaglig. En stor
\VUgras{viadukt} \textsuperscript{(16)} byggdes för att möjliggöra en järnväg;
\VUgras{kuranstaltens} \VUgras{park} \textsuperscript{(18)}, i vilken man ger
\VUgras{konserter}, är anlagd på ett förtjusande sätt. Flera prydliga
« \VUgras{villor} » \textsuperscript{(19)} byggdes runt kasinot \textsuperscript{(21)}, och en
mycket väl underhållen \VUgras{landsväg} \textsuperscript{(22)} underlättar förbindelserna
mycket.

%%K t09-c1-10-1 p
10. --- En enkel \VUgras{slänt} \textsuperscript{(23)} skiljer \VUgras{vägen}
\textsuperscript{(20)} från den egentliga \VUgras{dalen} \textsuperscript{(25)}, i vars botten
en prydlig liten \VUgras{sjö} \textsuperscript{(26)} ligger som föder en klar \VUgras{bäck}
\textsuperscript{(24)}.

%%K t09-c1-11-1 p
11. --- På andra sidan dalen bearbetas en \VUgras{järngruva} \textsuperscript{(28)}. Flera
gånger gick jag för att se \VUgras{gruvarbetarna} stiga ner i \VUgras{schaktet}, och till och
med gick jag in i \VUgras{gången} i vilken de tillbringar sin dag, medan de bryter
\VUgras{malmen} som innehåller
\VUgras{metallen}.

%%K t09-c1-12-1 p
12. --- Ett viktigt \VUgras{järnbruk} byggdes vid gruvan för att genast utnyttja de rikedomar
som bryts ur den. \VUgras{Masugnen} \textsuperscript{(29)}, upphettad med det \VUgras{stenkol}
som är rikligt här, gör det möjligt att snabbt och billigt få \VUgras{tackjärn}.

%%K t09-c1-13-1 p
13. --- Inte långt från gruvan gräver man ett \VUgras{stenbrott} \textsuperscript{(30)}.
Varken \VUgras{granit}, eller \VUgras{porfyr}, eller ens
\VUgras{sandsten} hugger den gamle \VUgras{stenhuggaren} med sin \VUgras{hacka},
utan enkel \VUgras{huggen sten} för husbygge.

%%K t09-c1-14-1 p
14. --- En av det landets vackraste prydnader är \VUgras{granarna} \textsuperscript{(31)}.
Överallt i botten av ett \VUgras{ravin} \textsuperscript{(32)}, vid stranden av en
\VUgras{fors} \textsuperscript{(33)}, av en \VUgras{avgrund} \textsuperscript{(34)} eller av
ett stup, skiftar deras tunna barr i grönt.

%%K t09-tit-6 sub
\VUsaut{3.0mm}
\VUpk{10.2pt}{40}{Att gå till fots.}
\VUsaut{3.0mm}

%%K t09-c1-15-1 p
15. --- Jag begagnar mitt lov för att gå \VUgras{långa} promenader.
\VUgras{Vägarna} \textsuperscript{(35)} som slingrar sig uppför berget gör det
möjligt att, utan att bry sig om avståndet, tillryggalägga flera
\VUgras{kilometer} och ofta flera \VUgras{mil}.

%%K t09-c1-15-2 p
\VUgras{Gränsstenar} \textsuperscript{(36)} och \VUgras{vägvisare}
\textsuperscript{(37)} märker ut vägarna på vilka man ofta möter en ung
\VUgras{bergsbo} \textsuperscript{(38)} med en \VUgras{ränsel}
\textsuperscript{(40)} vid sidan, som bär ett \VUgras{murmeldjur} \textsuperscript{(39)},
eller någon \VUgras{zigenare} \textsuperscript{(41)} vars \VUgras{pälsmössa}
\textsuperscript{(42)} egendomligt kontrasterar mot den gamla trasiga \VUgras{kapuschongen}
\textsuperscript{(43)}. Han leder i en
\VUgras{kedja en björn} \textsuperscript{(44)} som är dresserad.

%%K t09-tit-7 sub
\VUpk{10.2pt}{40}{Bergsbestigning.}
\VUsaut{3.0mm}

%%K t09-c1-16-1 p
16. --- Nästan varje dag går jag med en \VUgras{förare} \textsuperscript{(48)} för att göra en
\VUgras{bestigning}, och jag är mycket stolt när jag, med en
\VUgras{ränsel} \textsuperscript{(47)} på ryggen, och « \VUgras{alpstaven} » i
handen, som en riktig \VUgras{turist} \textsuperscript{(46)} stormar
\VUgras{klipporna} \textsuperscript{(45)}.

%%K t09-c1-17-1 p
17. --- Från ett bergs topp verkar slättens invånare inte större än myror; en
\VUgras{fårhjord} \textsuperscript{(50)} som betar på en \VUgras{betesmark}
\textsuperscript{(49)} verkar en barnleksak. En \VUgras{tall} \textsuperscript{(51)} eller
också ett \VUgras{litet trähus} \textsuperscript{(52)} ter sig knappt som en fläck på
grönskan.

%%K t09-c1-18-1 p
18. --- Under de utflykterna möter vi ofta på \VUgras{stigen} \textsuperscript{(53)} en
\VUgras{gemsjägare} \textsuperscript{(54)} som på sin skuldra bär det djur han just har dödat.

%%K t09-c1-19-1 p
19. --- Jag lade i mitt herbarium en mycket märkvärdig gren \VUgras{ormbunke}
\textsuperscript{(55)} och en vacker rosig kvist \VUgras{ljung} \textsuperscript{(57)} som jag
plockade vid ingången till en liten
\VUgras{grotta} \textsuperscript{(56)} under min sista promenad.

%%K t09-tit-8 sub
\VUsaut{3.0mm}
\VUcentre{12.0pt}{0}{\textit{Andra scenen.}}
\VUsaut{3.0mm}

%%K t09-tit-9 sub
\VUpk{11.4pt}{40}{Skogen. --- Jakten.}
\VUsaut{3.0mm}

%%K t09-c2-01-1 p
1. --- I går tillbringade jag en härlig dag i skogen. Den flit som råder bland smådjuren och
\VUgras{insekterna} \textsuperscript{(5)} av vilka den bebos intresserade mig mycket.

%%K t09-c2-01-2 p
\VUgras{Spindeln} \textsuperscript{(1)} som väver sitt \VUgras{nät}
\textsuperscript{(2)} mellan två grenar för att fånga \VUgras{flugan} \textsuperscript{(3)}
som far förbi, \VUgras{snigeln} \textsuperscript{(4)} som mödosamt bär sitt hus, ekoxen som
söker efter sitt byte, den flitiga myran som mycket ivrigt håller på vid \VUgras{myrstacken}
\textsuperscript{(6)}, alla de små gick, kom, arbetade med utomordentlig flit.

%%K t09-c2-02-1 p
2. --- En \VUgras{orm} \textsuperscript{(7)} som jag vid första anblicken höll för en
\VUgras{huggorm}, men som inte var något annat än en vanlig
\VUgras{snok}, hade gömt sig under en trädstam, och stack då och då ut sitt tunna
lilla huvud som alltid tycktes berett att bita. Framför mig kröp en tjock
\VUgras{snigel} \textsuperscript{(8)} tungt fram sedan den hade slukat ett av de
välsmakande \VUgras{smultronen} \textsuperscript{(11)} som växer rikligt där.

%%K t09-c2-03-1 p
3. --- Marken var mattbelagd med \VUgras{skogsblommor}; \VUgras{gullvivor}
\textsuperscript{(9)}, \VUgras{vintergröna} \textsuperscript{(10)}, och många andra växter som
jag inte kände, blommade mellan \VUgras{grästuvor} \textsuperscript{(12)}.

%%K t09-c2-04-1 p
4. --- I den trakten är \VUgras{tryffeln} nästan lika vanlig som
\VUgras{svampen} \textsuperscript{(14)} och en \VUgras{griskulting}
\textsuperscript{(13)} som tillhörde en skogvaktare, sökte glupskt efter om den skulle lyckas
hitta några.

%%K t09-tit-10 sub
\VUsaut{3.0mm}
\VUpk{10.2pt}{40}{Tjuvjakt.}
\VUsaut{3.0mm}

%%K t09-c2-05-1 p
5. --- Jag bevistade \VUgras{slutet} av en scen som roade mig mycket. Med hjälp av luktsinnet
hos sin \VUgras{taxhund} \textsuperscript{(15)}, hade
\VUgras{skogvaktaren} \textsuperscript{(16)} just överraskat en
\VUgras{tjuvskytt} \textsuperscript{(22)}, i det ögonblick då denne gjorde sig
redo att bära bort \VUgras{villebrådet} \textsuperscript{(17)} som han hade dödat. Eftersom
han hellre övergav sitt byte än lät sig gripas, hade tjuvskytten flytt, och en \VUgras{hare}
\textsuperscript{(18)}, en \VUgras{hind} \textsuperscript{(19)}, ett \VUgras{rådjur}
\textsuperscript{(20)} och ett
\VUgras{vildsvin} \textsuperscript{(21)} låg på gräset till skogvaktarens glädje.

%%K t09-c2-06-1 p
6. --- Mångfalden av de träd som skogen innehåller är häpnadsväckande.
\VUgras{Bokarna} \textsuperscript{(23)} och \VUgras{björkarna}
\textsuperscript{(26)}, \VUgras{tallarna}, som ger \VUgras{kådan} \textsuperscript{(27)},
\VUgras{ekarna} \textsuperscript{(29)} och också många andra blandar sina grenverk.

%%K t09-c2-06-2 p
\VUgras{Murgrönan} \textsuperscript{(24)}, som klänger sig fast vid de höga
trädens stam, färgar \VUgras{skogen} \textsuperscript{(25)} med ett glänsande grönt, som
behagligt förenar sig med den ljusare och blekgröna färgtonen hos
\VUgras{mossan} \textsuperscript{(31)} i vilken \VUgras{gröngölingen}
\textsuperscript{(30)} söker insekter.

%%K t09-tit-11 sub
\VUsaut{3.0mm}
\VUpk{10.2pt}{40}{Skogslandskap.}
\VUsaut{3.0mm}

%%K t09-c2-07-1 p
7. --- De gamla ekarna hänförde mig. Deras tjocka vridna \VUgras{rötter}
\textsuperscript{(32)}, till hälften uppkomna ur marken, ger dem ett praktfullt utseende av
kraft och styrka, och jag frågar mig hur \VUgras{ägaren} \textsuperscript{(35)} kan förmå sig
att låta fälla dem och utlämna dem åt
\VUgras{sågen} \textsuperscript{(34)} hos \VUgras{vedhuggaren}
\textsuperscript{(33)}. De stackars \VUgras{stammarna} \textsuperscript{(36)}, berövade sin
\VUgras{bark} \textsuperscript{(37)} eller kluvna av
\VUgras{yxan} \textsuperscript{(38)} hos \VUgras{dagsverkaren}
\textsuperscript{(39)} bedrövar mig.

%%K t09-c2-08-1 p
8. --- Nära intill hade en gammal \VUgras{kolare} \textsuperscript{(41)} med sin
\VUgras{skovel gjort i ordning en hög} \textsuperscript{(42)} kol, och en gammal
kvinna bar en \VUgras{knippa} \textsuperscript{(40)} \VUgras{torr ved} av kvistar från de
fällda träden.

%%K t09-tit-12 sub
\VUsaut{3.0mm}
\VUpk{10.2pt}{40}{Jakt.}
\VUsaut{3.0mm}

%%K t09-c2-09-1 p
9. --- Jag ville gå för att vila några ögonblick i \VUgras{skogvaktarens hus}
\textsuperscript{(46)} men, när jag kom fram till den \VUgras{lilla sjön}
\textsuperscript{(43)}, på vilken en vacker \VUgras{svan} \textsuperscript{(45)} simmar,
märkte jag, att en nyligen inträffad \VUgras{översvämning} \textsuperscript{(44)} hade
förvandlat vägen till ett verkligt \VUgras{träsk}. Jag måste vika av, och, medan jag gick
förbi \VUgras{cedern} \textsuperscript{(49)}, \VUgras{popplarna} \textsuperscript{(48)}, och
\VUgras{almen} \textsuperscript{(47)} som står vid skogens rand, såg jag komma
fram ur ett \VUgras{snår} \textsuperscript{(54)} en praktfull \VUgras{hjort}
\textsuperscript{(50)}, förföljd av ett envist \VUgras{hundkoppel} \textsuperscript{(51)} som
också eggades av \VUgras{hornet} \textsuperscript{(53)} hos \VUgras{jägarna till häst}
\textsuperscript{(52)}. Det stackars djuret var alldeles uttömt. Det kom och föll döende några
steg från mig.

%%K t09-c2-10-1 p
10. --- Skogvaktarens son, som hade lockats dit av \VUgras{drevjaktens} buller, kom ut ur
huset och vi gick tillsammans in i skogen igen. Han hade sett en
\VUgras{skata} \textsuperscript{(56)} på en gren av en gammal ek. Han menade, att
hennes bo troligen är gömt i \VUgras{lövverket} \textsuperscript{(58)} och han ville ta boet.

%%K t09-c2-10-2 p
Vig som en \VUgras{ekorre} \textsuperscript{(61)}, klättrade han från
\VUgras{gren} \textsuperscript{(55)} till gren, men han fann ingenting och lyckades
bara skrämma upp en gammal \VUgras{kråka} \textsuperscript{(60)} som satt på en
\VUgras{tjock gren} \textsuperscript{(57)}.
\VUsaut{4.0mm}
\VUfilet{22mm}
